\section{Motivating Descriptives}

Two descriptive patterns motivate the empirical analysis. First, ABAWD waiver exposure is neither static nor purely statewide: over time, states shift from broad statewide waivers to narrower county carveouts, and within waived states many counties remain untreated in some years. Second, dollar store growth substantially outpaces other SNAP-authorized retail formats over the same period and is geographically widespread. Taken together, these patterns suggest that both the policy environment and the retail landscape vary meaningfully across counties and over time.

\subsection{ABAWD Waivers}

As discussed above, FNS allowed states to request waivers for geographic units other than individual counties, including labor market areas, metropolitan divisions, towns, and reservations. This flexibility mattered because it allowed states to construct waiver areas around broader labor markets rather than strictly around county-by-county eligibility. Pennsylvania's use of the Philadelphia Metropolitan Division is illustrative: that waiver geography bundled Bucks, Chester, Delaware, Montgomery, and Philadelphia Counties into a single labor-market area. In settings like this, county treatment can reflect inclusion in a broader administrative waiver area rather than a stand-alone county determination. For our purposes, this means the relevant variation is not only whether a state used a waiver, but also whether it implemented the waiver statewide or through partial carveouts.

Figure \ref{fig:Waivergrowth} plots the number of counties covered by explicit county-level waivers over time, split by rural and urban status. A first takeaway is that county-level waiver use was highly time-varying. The first plotted year, 2014, is a low point with only 247 county-waived counties. County-level waiver coverage then rises sharply to 1,873 counties in 2015 before gradually receding to 920 by 2019. This figure should be read together with the state-level scope measure below: because Figure \ref{fig:Waivergrowth} counts explicit county-level waivers only, the unusually low 2014 count does not imply that waiver use disappeared. Instead, it marks a transition away from the near-ubiquitous statewide coverage of the immediate post-recession period and toward a regime in which fewer areas were waived and a larger share of waiver activity occurred through more targeted carveouts. A second takeaway is that rural counties consistently account for the majority of county-waived areas, representing roughly 60 to 66 percent of county-waived counties in each year from 2014 to 2019.

\begin{figure}
    \centering
    \includegraphics[width = 0.8\textwidth]{2026_AAEA_Submission/Figures/3_1_0_0_county_conferral_growth_rural_share.jpeg}
    \caption{County-level ABAWD waivers by year}
    \label{fig:Waivergrowth}
\end{figure}

Figure \ref{fig:annual_waiver_status_counts} clarifies how the scope of waiver use changed across states. In 2012 and 2013, 35 and 34 states in the county analysis panel coverage, respectively, operated statewide waivers, while only four states in each year used county waivers. By 2014, statewide waivers had nearly disappeared, with only three statewide-waiver states and six county-waiver states. Waiver activity broadened again in 2015, but in a different form: by 2016, 22 states used county waivers while only 8 retained statewide waivers, and by 2019 county waivers remained present in 21 states while statewide waivers had fallen to just 2. The central descriptive fact is therefore not a monotone decline in waivers, but a re-sorting from broad statewide relief toward more spatially selective waiver administration. That shift is important for the empirical design because it creates within-state county-level variation in treatment exposure.

\begin{figure}
    \centering
    \includegraphics[width = 0.8\textwidth]{2026_AAEA_Submission/Figures/3_1_0_5_annual_state_waiver_status_counts.jpeg}
    \caption{Annual state ABAWD waiver status}
    \label{fig:annual_waiver_status_counts}
\end{figure}

Figure \ref{fig:Waivercountystatus} maps counties that were ever assigned a county-level waiver during the analysis period. County-level waiver exposure is geographically broad rather than concentrated in a small number of isolated distressed places. In our county panel, 64.4 percent of counties are ever county-waived. The map shows especially broad county-level exposure across much of the South and West, alongside more persistent untreated areas in parts of the Great Plains, Texas, and northern New England. This wide spatial footprint is useful empirically because it indicates that county-level waiver treatment is not confined to one narrow region of the country.

\begin{figure}
    \centering
    \includegraphics[width = 0.8\textwidth]{2026_AAEA_Submission/Figures/3_1_0_2_ever_waived_county_map.jpeg}
    \caption{Counties ever assigned a county-level ABAWD waiver}
    \label{fig:Waivercountystatus}
\end{figure}

Figure \ref{fig:Waivercountystatusmap} turns from ever-treated status to treatment timing and within-state scope. The figure colors both statewide and county-level waiver exposure, but isolates directly county-waived counties within each treated state-year. In 2014, only 179 counties are directly county-waived, compared with 578 other counties located in states with some waiver activity; equivalently, only 23.6 percent of counties in waiver-active states are directly county-waived. By 2016 that share rises to 52.3 percent, and by 2019 it remains elevated at 45.4 percent. The maps therefore show that treatment timing is heterogeneous not only across states, but also within states. Some states move from no waiver to partial waiver coverage, others move from statewide waivers to more selective county carveouts, and others lose waiver coverage entirely. This staggered and spatially incomplete rollout is precisely the source of identifying variation exploited in the reduced-form analysis.

\begin{figure}
    \centering
    \includegraphics[width = 0.8\textwidth]{2026_AAEA_Submission/Figures/3_1_0_6_county_waiver_county_map.jpeg}
    \caption{ABAWD waiver geography in selected analysis years}
    \label{fig:Waivercountystatusmap}
\end{figure}

\subsection{Retailer Formats}

Figure \ref{fig:retail_format_growth} uses SNAP retailer location data to trace county-average retailer stock by format, normalized to 2010 = 100. Dollar stores exhibit the strongest sustained expansion over the decade: their county-average stock index rises from 100 in 2010 to 160 by 2019, implying 60 percent growth. Convenience stores also expand, but much more modestly, ending the period at an index value of roughly 125. In contrast, multi-category retailers rise only to about 117, while supermarkets are essentially flat to slightly declining and end the decade below their 2010 level. Notably, the dollar store series overtakes convenience stores by the middle of the decade and then continues to separate from the rest of the retail distribution. The descriptive implication is that growth in SNAP-authorized retail access during the 2010s was disproportionately driven by dollar stores rather than by supermarkets or big-box formats.

\begin{figure}
    \centering
    \includegraphics[width = 0.8\textwidth]{2026_AAEA_Submission/Figures/3_1_1_0_retailer_format_stock_index.jpeg}
    \caption{SNAP retailer format growth by type}
    \label{fig:retail_format_growth}
\end{figure}

Figure \ref{fig:ds_stock_change} shows that dollar store expansion was geographically broad rather than confined to a few local hotspots. Among counties with at least one dollar store in 2010, the median county increased dollar store stock by 50 percent by 2019, the mean increase was 72.3 percent, 81.3 percent of counties experienced positive growth, and 31 percent at least doubled their stock. At the same time, 16.5 percent of counties had no dollar stores in 2010, which is why those counties appear in gray and do not contribute to the displayed percent-change distribution. Even with that caveat, the map makes clear that the expansion of dollar stores was a national phenomenon with substantial cross-county heterogeneity in intensity. That broad but uneven expansion is what makes it plausible that exposure to dollar store growth could interact meaningfully with localized changes in SNAP eligibility policy.

\begin{figure}[H]
    \centering
    \includegraphics[width = 0.8\textwidth]{2026_AAEA_Submission/Figures/3_1_1_7_ds_stock_change_map_2010_2019.jpeg}
    \caption{Dollar store stock growth by county, 2010 to 2019}
    \label{fig:ds_stock_change}
\end{figure}
